% =========================================================
\chapter{Lecture 3 : Access control basics and DAC}
% =========================================================

\begin{rembox}[title=In a nutshell]
\textbf{What is this lecture about?} \textbf{Access control} answers one question:
``is \emph{this} user allowed to do \emph{this} action on \emph{this} file?''. In
\textbf{DAC}, the \textbf{owner} of a file decides who can use it (like on your laptop or
on Facebook). We see how to store these permissions (matrix, ACLs, capabilities, roles,
groups), a classic trap (the \emph{confused deputy}), and how Linux and Windows do it in
practice. \\
\textbf{Why is it useful?} Access control is the \textbf{first line of defense} and is
everywhere: apps, databases, operating systems, hardware. Understanding it lets you read
and set file permissions correctly (\texttt{ls -l}, \texttt{chmod}, \texttt{sudo}) and
spot programs that can be tricked into using their rights for an attacker.
\end{rembox}

\section{What is access control?}

\begin{defnbox}[title=Access control]
Security mechanism that ensures that ``all accesses and actions on objects by principals
are \textbf{within} the security policy''.
\end{defnbox}

\begin{exbox}[title=Questions an access control system answers]
Can Alice read file \texttt{/users/Bob/readme.txt}? Can Bob open a TCP socket to
\texttt{http://www.abc.com/}? Can Charlie write to row 15 of the table GRADES? \\
The answer is either \textbf{authorized} (``has permission'') or
\textbf{unauthorized} (``access denied''): only events within the security policy may
happen.
\end{exbox}

\begin{rembox}[title=Why is access control so important?]
Access control is the \textbf{first line of defense}, thus \textbf{it is used
everywhere}:
\begin{itemize}
  \item \textbf{Applications} --- online social networks, email servers, cloud storage;
  \item \textbf{Middleware} --- database management systems (DBMS);
  \item \textbf{Operating system} --- access to files, directories, ports, \ldots;
  \item \textbf{Hardware} --- memory, registers, privileges.
\end{itemize}
\end{rembox}

\begin{defnbox}[title=Authentication vs.\ authorization]
\begin{itemize}
  \item \textbf{Authentication} --- the system needs to \textbf{bind the actor to a
        principal} before authorization. A \textbf{principal} is an abstract entity
        that is authorized to act (users, connections, processes).
  \item \textbf{Authorization} --- the system needs to decide whether the
        \textbf{principal is authorized}: this is \textbf{access control}.
\end{itemize}
\end{defnbox}

\begin{warnbox}[title=TCB]
The mechanisms that do authentication and authorization are in the \textbf{Trusted
Computing Base}!
\end{warnbox}

\section{Implementing access control}

\begin{propbox}[title={``Checks soup'' vs.\ reference monitor}]
\begin{itemize}
  \item \textbf{What NOT to do: ``checks soup''} --- add checks all over the program,
        implementing the decision \textbf{in-line} based on the policy
        (\texttt{if user == Alice or user == Bob: open(\ldots)}).
  \item \textbf{What you SHOULD do: systematic calls to a ``reference monitor''} ---
        all over the program add checks that \textbf{call the monitor}: they state the
        authorization required and provide evidence about principals and objects; a
        \textbf{``central'' subsystem} establishes whether the checks pass or not
        (e.g.\ Apache Shiro: \texttt{subject.isPermitted("user:delete:jsmith")}).
\end{itemize}
\end{propbox}

\begin{rembox}[title=Least common mechanism?]
A central reference monitor is a mechanism shared by everybody --- but this is fine: least
common mechanism is about the isolation of data, and a single well-tested monitor is
easier to verify (economy of mechanism, complete mediation).
\end{rembox}

\section{DAC vs.\ MAC}

\begin{defnbox}[title=Who decides the access policy? Two approaches]
\begin{itemize}
  \item \textbf{Discretionary Access Control (DAC)} --- \textbf{object owners} assign
        permissions (ownership of resources). E.g.\ Windows, Linux, macOS, social
        networks.
  \item \textbf{Mandatory Access Control (MAC)} --- a \textbf{central security policy}
        assigns permissions. For organizations with need for central controls:
        military (confidentiality), hospital (confidentiality and integrity), banking
        (integrity).
\end{itemize}
\end{defnbox}

\begin{rembox}[title={Origin of protection mechanisms (Lampson, 1971)}]
The original motivation was to keep \textbf{one user's malice or error from harming other
users}: by destroying/modifying another user's data, by reading/copying it without
permission, or by degrading the service another user gets (the extreme case being a
crash of the system). All these reasons are just as strong if ``user'' is replaced by
``\textbf{program}''.
\end{rembox}

\section{The access control matrix}

\begin{defnbox}[title=Access control matrix]
An \textbf{abstract} representation of all \textbf{permitted} triplets
\textbf{(subject, object, access right)} within a system. DAC policies are often
conceptualized as an access control matrix (Lampson).
\begin{itemize}
  \item \textbf{Subjects (principals)} --- entities within an IT system (a user, a
        process, a service).
  \item \textbf{Objects (assets)} --- resources that (some) subject may access or use
        (a file, a folder, a row in a database, the system's memory, a machine in the
        network, a printer, a page in a website).
  \item \textbf{Operations} --- in abstract, subjects can \textbf{observe} and/or
        \textbf{alter} objects (read, write, append, execute).
\end{itemize}
\end{defnbox}

\begin{exbox}[title=Access control matrix]
$S = \{\text{Alice}, \text{Bob}\}$, $O = \{\text{file1}, \text{file2},
\text{file3}\}$, $A = \{\text{read}, \text{write}\}$.
\begin{center}
\begin{tabular}{l|c|c|c}
      & file1      & file2      & file3      \\ \hline
Alice & read, write &            & read       \\
Bob   &            & read, write & read, write
\end{tabular}
\end{center}
Can Alice read file1? Can Bob write file1? Can file3 be written by Alice?
\end{exbox}

\begin{warnbox}[title=The matrix is an abstract concept]
Not for direct implementation: size $O(f \cdot u)$ ($f$ files, $u$ users), which is huge
with thousands of files or hundreds of users. Moreover it is:
\begin{itemize}
  \item \textbf{usually very sparse} --- wasteful;
  \item \textbf{inefficient to access} --- need to keep the matrix in fast memory?
  \item \textbf{not extensible} --- add a file? add a user?
\end{itemize}
\end{warnbox}

\section{Access Control Lists (ACLs)}

\begin{defnbox}[title=Access Control List (ACL)]
Associate permissions to \textbf{objects}: each object stores the \textbf{column} of the
matrix, i.e.\ the list of (subject, rights) pairs. \textbf{Blanks are not stored}.
\begin{center}
\texttt{file3: \{(Alice, read), (Bob, read/write)\}}
\end{center}
\end{defnbox}

\begin{propbox}[title=ACLs: pros and cons]
\textbf{Pros}: can be stored close to / with the resource; easy to determine who can
access a resource; easy to revoke rights by resource. \\
\textbf{Cons}: difficult to check at runtime; difficult to audit all rights of a user;
difficult to remove all permissions from a user (better remove authentication!);
difficult delegation.
\end{propbox}

\subsection{Role Based Access Control (RBAC)}

\begin{defnbox}[title=Role Based Access Control (RBAC)]
Systems have too many subjects that come and go (large dynamic ACLs), but subjects are
similar to each other and need the same rights (a doctor has the same privileges as
another doctor).
\begin{enumerate}
  \item assign permissions to \textbf{roles};
  \item assign \textbf{roles} to subjects;
  \item subjects select an \textbf{active role} --- they have the permissions of the
        active role.
\end{enumerate}
\end{defnbox}

\begin{warnbox}[title=RBAC problems]
\begin{enumerate}
  \item \textbf{Role explosion} --- temptation to create fine-grained roles, denying the
        benefits of RBAC.
  \item \textbf{Limited expressiveness} of simple RBAC --- problems implementing least
        privilege; some roles are relative (``Alice's doctor'' vs.\ ``any doctor'').
  \item \textbf{Separation of duty is difficult} --- ``two doctors are needed to
        authorize a procedure'': the mechanism needs to ensure they are distinct!
\end{enumerate}
\end{warnbox}

\subsection{Group Based Access Control}

\begin{defnbox}[title=Group Based Access Control]
Observation: some permissions are always needed together (e.g.\ access to sockets and
network interface).
\begin{enumerate}
  \item assign permissions to access objects to \textbf{groups};
  \item assign subjects to \textbf{groups};
  \item subjects have the permissions of \textbf{all} their groups.
\end{enumerate}
\textbf{Negative permissions} (explicit denials for a subject) can be added to implement
fine-grained policies.
\end{defnbox}

\begin{rembox}[title=RBAC vs.\ groups]
With \textbf{roles}, a subject has the permissions of the \emph{one active} role; with
\textbf{groups}, a subject has the union of the permissions of \emph{all} its groups.
\end{rembox}

\begin{propbox}[title=Negative permissions are checked first]
If group permissions were checked first, a group would grant the access and the
negative permission would never be applied. Checking \textbf{negative permissions first}
is the fail-safe choice.
\end{propbox}

\section{Capabilities}

\begin{defnbox}[title=Capabilities]
Associate permissions to \textbf{subjects}: each subject stores the \textbf{row} of the
matrix, i.e.\ the list of (object, rights) pairs. \textbf{Blanks are not stored}.
\begin{center}
\texttt{Alice: \{(file1, read/write), (file3, read)\}}
\end{center}
\end{defnbox}

\begin{propbox}[title=Capabilities: pros and cons]
\textbf{Pros}: can be stored with the subject (portable!); easy to audit all subject
permissions; delegating is ``simple''. \\
\textbf{Cons}: revoking permission on one object is hard; transferability --- once the
capability is given, how can we prevent sharing?; authenticity --- how to check it?
\end{propbox}

\begin{keybox}[title=ACL vs.\ capabilities]
\begin{center}
\begin{tabular}{l|l|l}
                       & \textbf{ACL}           & \textbf{Capabilities} \\ \hline
Stored with            & the object (column)    & the subject (row)     \\
Who can access $o$?    & easy                   & hard                  \\
What can $s$ access?   & hard                   & easy                  \\
Revocation             & easy per object        & hard per object       \\
Delegation             & difficult              & ``simple''
\end{tabular}
\end{center}
\end{keybox}

\begin{rembox}[title=History of capability-based security]
Introduced by Dennis \& Van Horn (1966); used in the Plessey System 250 (1972), CAP
computer (1976), Intel iAPX 432 (1981). The industry then adopted UNIX and its ACLs
(1982, Linux 1991). Revival: seL4 formally verified microkernel (2009), CHERI (2015),
Google Fuchsia (2016), WASI (2019), ARM Morello (2022), Scala capabilities at EPFL.
\end{rembox}

\section{Ambient authority and the confused deputy}

\begin{defnbox}[title=Ambient authority]
\textbf{Ambient authority} is used by a subject if, for an action to succeed, it
\textbf{only} needs to specify the \textbf{operation} and the \textbf{names} of the
involved object(s). The \textbf{subject} (with authority) is \textbf{implicit}:
in \texttt{open("file1", "rw")} the subject is missing, it is understood to be the
process owner.
\end{defnbox}

\begin{propbox}[title=Ambient authority: pros and cons]
\textbf{Pro}: no need to repeat the subject all the time (usability). \\
\textbf{Cons}: the program cannot check permissions for whom it acts; least privilege
becomes harder to enforce; \textbf{confused deputy problem}!
\end{propbox}

\begin{defnbox}[title=Confused deputy problem]
Problem with ambient authority: a \textbf{privileged program} can be \textbf{tricked to
misuse its authority}.
\end{defnbox}

\begin{exbox}[title=Pay-per-use compiler]
The compiler receives \texttt{(input, output)}, compiles \texttt{input}, writes usage
into \texttt{bill} and errors into \texttt{output}.
\begin{center}
\begin{tabular}{l|c|c|c}
         & input & output      & bill        \\ \hline
Alice    & write & read        & read        \\
Compiler & read  & read, write & read, write
\end{tabular}
\end{center}
Alice cannot write \texttt{bill}, but she calls the compiler with
\texttt{(input, bill)}: the compiler writes the errors into \texttt{output = bill}
\textbf{with its own authority} $\Rightarrow$ \texttt{bill} is corrupted. The compiler
is the \textbf{confused deputy}. (Same idea: a student with a card asks someone
authorized to badge them into a room.)
\end{exbox}

\begin{algbox}[title=How to avoid confused deputies]
Ambient authority is used for convenience in many real systems (OS services, web
servers, \ldots). Solutions:
\begin{enumerate}
  \item re-implement access control in the privileged process;
  \item let the privileged process check authorization for Alice;
  \item \textbf{capabilities can help}! The compiler holds a capability for
        \texttt{bill}; to compile, Alice must give a capability for the
        \texttt{output} file --- she cannot give a capability for writing
        \texttt{bill}, so she cannot confuse anyone.
\end{enumerate}
\end{algbox}

\section{DAC in practice: UNIX}

\begin{defnbox}[title=UNIX principals and groups]
\begin{itemize}
  \item \textbf{User Identities (UIDs)} and \textbf{Group Identities (GIDs)}:
        originally 16-bit (now 32-bit) numbers. Special UIDs: $-2$, $0$, $1$, \ldots
  \item Each user has its own directory \texttt{/home/username}; user accounts are in
        \texttt{/etc/passwd} (\texttt{username:password:UID:GID:info:home:shell}).
  \item Users belong to one or more groups: primary group (\texttt{/etc/passwd}),
        other groups (\texttt{/etc/group}).
\end{itemize}
\end{defnbox}

\begin{propbox}[title=UNIX security architecture]
\begin{itemize}
  \item \textbf{Everything is a file}; each user ``owns'' a set of files.
  \item Each file has a simple \textbf{ACL} expressing its access control policy. The
        owner can \textbf{modify the permissions of the file} (but cannot give away
        ownership). System files are owned by special users that can make system
        operations.
  \item All processes run by a user run \textbf{with that user's privileges}
        $\Rightarrow$ \textbf{ambient authority}!
\end{itemize}
\end{propbox}

\begin{defnbox}[title=UNIX file ACLs]
Each file is assigned an \textbf{owner UID} and a \textbf{GID}, and has \textbf{9
permission bits}: 3 actions (\textbf{r}ead, \textbf{w}rite, e\textbf{x}ecute) $\times$ 3
subjects (\textbf{owner}, \textbf{group}, \textbf{other}). Plus 3 attributes:
\texttt{suid}, \texttt{sgid} and \texttt{sticky}. \\
Different semantics for \textbf{directories}: Read $\to$ list files, Write $\to$ add
files, Exec $\to$ traverse the directory (\texttt{cd}).
\end{defnbox}

\begin{exbox}[title=Reading \texttt{ls -l}]
\texttt{drwxrwxr-x 1 catronco catronco 4096 Sep 16 14:23 exampledir} \\
First char: \texttt{d} = directory, \texttt{-} = file; then \texttt{rwx} for owner,
\texttt{rwx} for group, \texttt{r-x} for others; then number of links, owner, group,
size, last modified, filename. \\
The owner changes permissions with \texttt{chmod}: \texttt{+r}, \texttt{-w},
octal \texttt{666}, \texttt{662}, sticky \texttt{+t} or \texttt{1666}, suid
\texttt{+s} or \texttt{4666}.
\end{exbox}

\begin{algbox}[title=UNIX access control in action]
Compare the \texttt{UID}/\texttt{GID} of the \textbf{process} trying to perform the
action with the state of the file (owner, group, mode bits). \textbf{Order matters}:
\begin{enumerate}
  \item if the \texttt{UID} says you are the owner: check the bits for \textbf{owner};
  \item if not owner, but your group is the file's group: check the bits for
        \textbf{group};
  \item otherwise check the bits for \textbf{other}.
\end{enumerate}
The \texttt{root} user is \textbf{never denied access}.
\end{algbox}

\begin{warnbox}[title=Order matters]
Only \emph{one} class of bits is checked: an owner with \texttt{---} for owner but
\texttt{rwx} for group is denied, even if they are in the group.
\end{warnbox}

\begin{defnbox}[title=Super user (\texttt{root})]
Special user account with \textbf{User ID 0}: accesses system files and special
operations, can access anything ((almost) all security checks disabled).
\textbf{\texttt{root} is in the TCB!} \\
\textbf{Never login as root}: use \texttt{sudo} (run one command as root) or
\texttt{su} (switch user).
\end{defnbox}

\begin{defnbox}[title=\texttt{suid} / \texttt{sgid}]
Setuid and setgid bits indicate that a file is \textbf{not run with the privileges of
the launcher}, but \textbf{with the privileges of the owner} user/group. Useful for
normal users to access system services that need system privileges, respecting least
privilege, e.g.\ changing a password:
\texttt{-rwsr-xr-x 1 root root \ldots /bin/passwd}.
\end{defnbox}

\begin{warnbox}[title=Setuid root programs are dangerous]
They are in the TCB: how do you know that a \texttt{suid} program does what it is meant
to do, and \emph{only} what it is meant to do? (e.g.\ a setuid-root \texttt{python}
gives root to anyone.)
\end{warnbox}

\begin{defnbox}[title=Sticky bit]
``Restricted deletion bit'' (\texttt{chmod +t}).
\begin{itemize}
  \item \textbf{Directories}: prevents unprivileged users from removing or renaming a
        file in the directory \textbf{unless they own the file} (e.g.\ \texttt{/tmp}:
        users can only edit their own files).
  \item \textbf{Files}: historically prevented a program from being moved from swap
        for fast load; Linux currently ignores the bit.
\end{itemize}
\end{defnbox}

\begin{defnbox}[title=Nobody]
Special user (\textbf{User ID $-2$}) that owns no files and belongs to no user. Safer
user to execute code you do not know (particularly obfuscated code): limits damages if it
misbehaves / gets compromised.
\end{defnbox}

\begin{exbox}[title={Exercise --- ACL in Linux}]
\texttt{drwxr-x--- 2 manager\_A core\_devs 4096 Sep 24 11:35 /project\_beta} \\
\textbf{manager\_A} is the directory owner, a member of \texttt{core\_devs}, and is
sudo; \textbf{developer\_B} is a member of \texttt{core\_devs}; \textbf{intern\_C} is
\textbf{not} a member of \texttt{core\_devs}. \\
For each user, determine which of \texttt{cd /project\_beta}, \texttt{ls /project\_beta}
and \texttt{touch /project\_beta/<file>} succeed or fail, and explain why based on the DAC
rules. \\[4pt]
Same state, with a new subdirectory and a file inside it: \\
\texttt{drwxr-xr-x 2 developer\_B core\_devs 4096 Sep 24 12:00 /project\_beta/Foo} \\
\texttt{-rw-r--r-- 1 developer\_B core\_devs 512 Sep 24 12:05 \ldots/Foo/intern\_project.txt}
\begin{enumerate}[label=Q\arabic*:]
  \item \texttt{/project\_beta} is owned by manager\_A, but \texttt{Foo} is owned by
        developer\_B. Find a scenario (sequence of commands) that would have led to this
        situation.
  \item intern\_C runs \texttt{cd /project\_beta/Foo} then
        \texttt{cat intern\_project.txt}. Do they work?
  \item intern\_C runs \texttt{cat /project\_beta/Foo/intern\_project.txt} from their
        home directory. Does it work?
  \item Who do they need to contact to fix the problem?
\end{enumerate}
\end{exbox}

\section{DAC in practice: Windows}

\begin{defnbox}[title=Windows access control]
Principals = users, machines, groups, \ldots; objects = files, registry keys, printers,
\ldots
\begin{itemize}
  \item Each object has a \textbf{Discretionary Access Control List (DACL)}.
  \item Each process (or thread) has an \textbf{access token} with: the login user
        account (the process ``runs as'' this user), all groups of which the user is a
        member (recursively!), all privileges assigned to these groups.
\end{itemize}
The DACL is compared with the process' access token \textbf{when creating a handle to the
object}.
\end{defnbox}

\begin{defnbox}[title=DACL]
A list of \textbf{Access Control Entries (ACEs)}, each with a \textbf{type}
(negative/positive, i.e.\ access denied/allowed), a \textbf{principal}, and
\textbf{permissions} (more fine-grained than UNIX), plus flags and others. Negative
entries come \textbf{first} so that a denial cannot be bypassed by an earlier allow.
\end{defnbox}

\begin{exbox}[title=DACL evaluation]
ACE1: \textit{deny} Andrew read/write/execute; ACE2: \textit{allow} Group A write; ACE3:
\textit{allow} Everyone read/execute. Thread A (Andrew, groups A, B, C) is denied by ACE1.
Thread B (Jane, group A) is granted write (ACE2) and read/execute (ACE3).
\end{exbox}

\begin{rembox}[title=Least privilege by default]
Windows runs with least privilege by default; elevated rights require
``Run as administrator''.
\end{rembox}

\begin{keybox}[title=Summary]
\begin{itemize}
  \item \textbf{Discretionary Access Control}: owners establish permissions.
  \item Conceptualized as an \textbf{access control matrix}.
  \item Implemented in two flavors: \textbf{access control lists} (permissions
        associated to objects) and \textbf{capabilities} (permissions associated to
        subjects).
  \item When relying on access control, always think about \textbf{confused deputies}.
\end{itemize}
\end{keybox}
